\documentclass[9pt, a4paper, twocolumn]{article}
\usepackage[utf8]{inputenc}
\usepackage[margin=0.5in]{geometry}
\usepackage{amsmath, amssymb, amsfonts}
\usepackage{titlesec}
\usepackage{enumitem}
\usepackage{tikz}

\titleformat{\section}{\small\bfseries\uppercase}{}{0pt}{}[\titlerule]
\titleformat{\subsection}{\footnotesize\bfseries}{}{0pt}{}
\titlespacing{\section}{0pt}{6pt}{4pt}
\titlespacing{\subsection}{0pt}{4pt}{2pt}
\setlist[itemize]{noitemsep, topsep=0pt, leftmargin=10pt}

\begin{document}

\section{I. Fluid Kinematics}
A fluid element's motion is the sum of translation, linear deformation, rotation, and angular deformation.

\subsection{1. Linear Deformation and Dilatation}
Linear strain rate in $x$-direction: $\epsilon_{xx} = \frac{\partial u}{\partial x}$.
Volumetric dilatation rate (rate of change of volume per unit volume):
\begin{equation*}
    \frac{1}{\delta \forall} \frac{d(\delta \forall)}{dt} = \frac{\partial u}{\partial x} + \frac{\partial v}{\partial y} + \frac{\partial w}{\partial z} = \nabla \cdot \vec{V}
\end{equation*}
For incompressible flow: $\nabla \cdot \vec{V} = 0$.

\subsection{2. Rotation and Vorticity}
Angular velocity $\vec{\omega}$ is the average rate of rotation of two perpendicular lines:
\begin{equation*}
    \omega_z = \frac{1}{2} \left( \frac{\partial v}{\partial x} - \frac{\partial u}{\partial y} \right), \quad \vec{\omega} = \frac{1}{2} \nabla \times \vec{V}
\end{equation*}
Vorticity $\vec{\zeta}$ is twice the angular velocity: $\vec{\zeta} = 2\vec{\omega} = \nabla \times \vec{V}$.
In cylindrical coordinates:
\begin{equation*}
    \zeta_z = \frac{1}{r} \frac{\partial (r V_\theta)}{\partial r} - \frac{1}{r} \frac{\partial V_r}{\partial \theta}
\end{equation*}
\textbf{Irrotational Flow:} $\nabla \times \vec{V} = 0$.

\subsection{3. Angular Deformation}
Shear strain rate $\gamma_{xy}$ is the rate of change of the angle between two perpendicular lines:
\begin{equation*}
    \gamma_{xy} = \frac{\partial v}{\partial x} + \frac{\partial u}{\partial y}
\end{equation*}

\section{II. Conservation of Mass (Continuity)}
Derived using an infinitesimal Control Volume (CV).
\begin{itemize}
    \item \textbf{General (Vector):} $\frac{\partial \rho}{\partial t} + \nabla \cdot (\rho \vec{V}) = 0$
    \item \textbf{Cartesian:} $\frac{\partial \rho}{\partial t} + \frac{\partial (\rho u)}{\partial x} + \frac{\partial (\rho v)}{\partial y} + \frac{\partial (\rho w)}{\partial z} = 0$
    \item \textbf{Cylindrical:} $\frac{\partial \rho}{\partial t} + \frac{1}{r}\frac{\partial (r \rho V_r)}{\partial r} + \frac{1}{r}\frac{\partial (\rho V_\theta)}{\partial \theta} + \frac{\partial (\rho V_z)}{\partial z} = 0$
    \item \textbf{Incompressible:} $\frac{\partial u}{\partial x} + \frac{\partial v}{\partial y} + \frac{\partial w}{\partial z} = 0$
\end{itemize}

\section{III. The Stream Function ($\psi$)}
Defined for 2D, incompressible flow. Automatically satisfies continuity.
\begin{itemize}
    \item \textbf{Cartesian:} $u = \frac{\partial \psi}{\partial y}, \quad v = -\frac{\partial \psi}{\partial x}$
    \item \textbf{Cylindrical:} $V_r = \frac{1}{r} \frac{\partial \psi}{\partial \theta}, \quad V_\theta = -\frac{\partial \psi}{\partial r}$
\end{itemize}
\textbf{Properties:} 
1. $\psi = \text{const}$ defines a streamline. 
2. Volume flow rate per unit depth: $q = \psi_2 - \psi_1$. 
3. Relationship to Irrotationality: $\nabla^2 \psi = -\zeta_z$. Therefore, for \textbf{irrotational} flow, $\psi$ must satisfy Laplace's equation: $\nabla^2 \psi = 0$. (Rotational flow does not).
\section{IV. Conservation of Momentum}
\subsection{1. Material Derivative}
$\frac{D\vec{V}}{Dt} = \frac{\partial \vec{V}}{\partial t} + (\vec{V} \cdot \nabla)\vec{V}$. 
Acceleration includes local ($\partial/\partial t$) and convective parts.

\subsection{2. Euler's Equation (Inviscid Flow)}
Assuming $\mu = 0$ (frictionless):
\begin{equation*}
    \rho \vec{g} - \nabla p = \rho \frac{D\vec{V}}{Dt}
\end{equation*}
\textbf{In Cylindrical Coordinates:}
\begin{equation*}
    \rho \left( \frac{\partial V_r}{\partial t} + V_r \frac{\partial V_r}{\partial r} + \frac{V_\theta}{r} \frac{\partial V_r}{\partial \theta} - \frac{V_\theta^2}{r} + V_z \frac{\partial V_r}{\partial z} \right) = -\frac{\partial p}{\partial r} + \rho g_r
\end{equation*}
\begin{equation*}
    \rho \left( \frac{\partial V_\theta}{\partial t} + V_r \frac{\partial V_\theta}{\partial r} + \frac{V_\theta}{r} \frac{\partial V_\theta}{\partial \theta} + \frac{V_r V_\theta}{r} + V_z \frac{\partial V_\theta}{\partial z} \right) = -\frac{1}{r} \frac{\partial p}{\partial \theta} + \rho g_\theta
\end{equation*}
\begin{equation*}
    \rho \left( \frac{\partial V_z}{\partial t} + V_r \frac{\partial V_z}{\partial r} + \frac{V_\theta}{r} \frac{\partial V_z}{\partial \theta} + V_z \frac{\partial V_z}{\partial z} \right) = -\frac{\partial p}{\partial z} + \rho g_z
\end{equation*}

Bernoulli Equation is the integration of Euler's along a streamline for steady, incompressible flow.

\subsection{3. Navier-Stokes Equations (Viscous Flow)}
For Newtonian, incompressible fluids ($\tau_{ij} = \mu \gamma_{ij}$):
\begin{equation*}
    \rho \left( \frac{\partial \vec{V}}{\partial t} + (\vec{V} \cdot \nabla)\vec{V} \right) = -\nabla p + \rho \vec{g} + \mu \nabla^2 \vec{V}
\end{equation*}

\section{V. Potential Flow (Inviscid + Irrotational)}
Since $\nabla \times \vec{V} = 0$, there exists a velocity potential $\phi$: $\vec{V} = \nabla \phi$.
\begin{itemize}
    \item \textbf{Cartesian:} $u = \frac{\partial \phi}{\partial x}, \quad v = \frac{\partial \phi}{\partial y}$
    \item \textbf{Cylindrical:} $V_r = \frac{\partial \phi}{\partial r}, \quad V_\theta = \frac{1}{r}\frac{\partial \phi}{\partial \theta}$
\end{itemize}
Governing Eq: $\nabla^2 \phi = 0$ (Laplace). Cauchy-Riemann relations connect $\phi$ and $\psi$: $u = \frac{\partial \phi}{\partial x} = \frac{\partial \psi}{\partial y}$ and $v = \frac{\partial \phi}{\partial y} = -\frac{\partial \psi}{\partial x}$.

\subsection{1. Elementary Flows}
\begin{itemize}
    \item \textbf{Uniform Flow (U):} $\psi = Uy, \phi = Ux$
    \item \textbf{Source/Sink:} ($m$: volumetric flow rate per unit length). Units: $[m^2/s]$ $V_r = \frac{m}{2\pi r}, V_\theta = 0 \Rightarrow \psi = \frac{m}{2\pi}\theta, \phi = \frac{m}{2\pi}\ln r$
    \item \textbf{Irrotational Vortex ($\Gamma$):} $V_r = 0, V_\theta = \frac{\Gamma}{2\pi r} \Rightarrow \psi = -\frac{\Gamma}{2\pi}\ln r, \phi = \frac{\Gamma}{2\pi}\theta$
    \item \textbf{Doublet ($\kappa$):} $\psi = -\frac{\kappa \sin \theta}{r}, \phi = \frac{\kappa \cos \theta}{r}$
\end{itemize}

\subsection{2. Superposition of Flows}
\begin{itemize}
    \item \textbf{Half-Body:} Uniform + Source. Stagnation point at $r=m/(2\pi U), \theta = \pi$.
    \item \textbf{Rankine Oval:} Uniform + Source + Sink.
    \item \textbf{Flow Around Cylinder:} Uniform + Doublet. Surface is $r = R = \sqrt{\kappa/U}$.
    \item \textbf{Cylinder with Circulation:} Uniform + Doublet + Vortex. 
    Lift $L = \rho U \Gamma$ (Kutta-Joukowski).
\end{itemize}

\section{VI. Exact Solutions for Viscous Flow}
\subsection{1. Plane Poiseuille Flow (Fixed Parallel Plates)}
Assumptions: Steady, 2D ($v=w=0$), fully developed ($\partial u / \partial x = 0$).
Continuity $\Rightarrow u = u(y)$. Momentum $\Rightarrow \frac{dP}{dx} = \mu \frac{d^2 u}{dy^2} = \text{const}$.
Solution for gap $2h$ (plates at $y = \pm h$):
\begin{equation*}
    u(y) = \frac{1}{2\mu} \left(\frac{dP}{dx}\right) (y^2 - h^2)
\end{equation*}
$u_{max}$ at centerline. $V_{avg} = \frac{2}{3}u_{max}$.

\subsection{2. Hagen-Poiseuille Flow (Pipe)}
Steady, laminar, fully developed in pipe of radius $R$.
\begin{equation*}
    V_z(r) = \frac{1}{4\mu} \left( \frac{dP}{dz} \right) (r^2 - R^2), \quad Q = -\frac{\pi R^4}{8\mu} \frac{dP}{dz}
\end{equation*}

\subsection{3. Couette Flow}
Flow driven by a moving plate at $y=b$ with velocity $U$:
\begin{equation*}
    u(y) = U\frac{y}{b} + \frac{1}{2\mu}\left(\frac{dP}{dx}\right)(y^2 - by)
\end{equation*}

\section{VII. Summary of Assumptions}
\begin{itemize}
    \item \textbf{Inviscid:} Neglect $\mu$. Valid far from boundaries.
    \item \textbf{Irrotational:} $\nabla \times \vec{V} = 0$. Potential $\phi$ exists.
    \item \textbf{Incompressible:} $\rho = \text{const}$. $\nabla \cdot \vec{V} = 0$.
    \item \textbf{Fully Developed:} Velocity profile does not change in flow direction.
\end{itemize}

\end{document}